\section{Limitations\label{appendix: limitation}}
While \BenchName{} attempts to circumvent the current pitfalls of LLM benchmarks, some limitations remain:

\paragraph{Human-dependent Annotation} One limitation of \BenchName{} is its reliance on human-annotated subgoals to measure the progress rate of LLM agents. The subjectivity of annotators calls for multiple verifications just to ensure the uniformity of the data. Also, as the complexity of tasks increases, the number of subgoals and the granularity required in annotation can grow significantly. This makes the process less scalable and more labor-intensive. One alternative is to use LLMs instead of humans to annotate these subgoals, but all LLMs currently underperform on \BenchName{} tasks and cannot accurately generate subgoals. We look forward to better LLMs that could autonomize planning subgoal annotation.

\paragraph{Benchmarking Real-World Problems} The price to pay for a standardized and definitive agent benchmark is to evaluate agents in simulated environments. However, it's important to also benchmark on real-world problems for future applications. Currently, there are a few challenges with benchmarking on real-world problems that we hope to overcome in subsequent work: (1) Variable ground truth labels: most real-world environments are ever-changing, e.g., contents on a web page, and this could lead to changes in states and labels for LLM agents, creating difficulties for benchmarking. (2) Security measures: we benchmark not only on obtaining information from environments but also on operating within and altering these environments. In real-world scenarios, it could be dangerous to unleash an LLM agent, e.g., on the Internet, which could lead to harm or generate malicious information. Therefore, it is very important yet challenging to constrain the LLM agent's action space in the real world while still offering it the freedom to accomplish tasks.




\section{Ethics and Societal Impact\label{appendix: ethics}}
This paper presents work whose goal is to advance the field of Machine Learning. Potential societal implications include we allow LLM agents to access online API in tool-operation evaluation. LLM agents could access and edit online information including Google Sheet and To-do list. However, we made sure that no personal information is leaked and no generated content is distributed online during the evaluation process.